\section{Step 30 Clip the reduction-tile exclusive end}\label{matmul-s30}

At \texttt{block\_reduction}, form L+T in \texttt{a2} and clip it against K. This endpoint controls the number of products contributed in the current tile. A complete last tile is allowed; a shorter last tile must stop before reading an absent input pair.

\Needspace{10\baselineskip}
\subsection{Write and inspect}
Replace the temporary branch and nop in \texttt{t\_reduction\_bound}. Keep its surrounding labels and region markers.

\MARSCodeLabel{MIPS assembly}
\begin{minted}[samepage=true]{mars-idea}
    addu $a2, $t0, $s3
    sltu $t9, $s1, $a2
    beq $t9, $zero, reduction_bound
    nop
    move $a2, $s1
\end{minted}
Use \texttt{s30-1}, \texttt{s30-2} from the companion construction inputs. Pause at \texttt{matmul\_done} before the caller stores the returned count.

\subsection{Expected observation}
The first default reduction tile ends at two. For K=1,T=3, \texttt{a2} clips to one.

\Needspace{10\baselineskip}
\subsection{Check your understanding}
What bounds the last reduction pair?

\begin{itemize}
  \item \textbf{A} min(L+T,K).
  \item \textbf{B} min(J+T,N).
  \item \textbf{C} The size of the mapped data segment.
\end{itemize}
